% Lösungen Tangente, Normale, Schnittwinkel (zusammenbau v0.4, Heft)
\einheitenkopf{Tangente, Normale, Schnittwinkel}

\erg{43}{a) ein Punkt – ableiten und die Stelle einsetzen. \quad b) $f(2) = 4$, $f'(x) = 3x^2 - 2$, $m = f'(2) = 10$; $n = 4 - (10) \cdot (2) = -16$; $t(x) = 10x - 16$ \quad c) $f(1) = 2 - 1 = 1$ – $Q$ liegt auf dem Graphen; $m = f'(1) = 6 - 1 = 5$, $n = 1 - 5 = -4$; $t(x) = 5x - 4$}
\erg{44}{a) $f(3) = -2$, $f'(x) = 3x^2 - 6x$, $m = f'(3) = 9$; $n = -2 - (9) \cdot (3) = -29$; $t(x) = 9x - 29$ \quad b) $f(0) = 2 - 3 = -1$; $f'(x) = 2e^x$, $m = f'(0) = 2e^0 = 2$ (nicht $0$); $t(x) = 2x - 1$ \quad c) $m = k'(15) = -2e^{-1{,}5} \approx -0{,}45$; $n = k(15) - 15m \approx 43{,}39 + 6{,}69 \approx 50{,}08$; $g(x) \approx -0{,}45x + 50{,}08$; $g(x) \le 30$: beim Teilen durch die negative Steigung dreht sich das Zeichen, Grenze $x = 45$ – ab Tag $45$ \quad d) $f(0) = 4$; $f'(x) = -3\mathrm{cos}(x)$, $m = f'(0) = -3$; $t(x) = -3x + 4$ \quad e) $(-1 + \sqrt{6})^2 = 7 - 2\sqrt{6}$; $f'(-1 + \sqrt{6}) = (7 - 2\sqrt{6} - 2 + 2\sqrt{6} - 5) \cdot e^{-1 + \sqrt{6}} = 0$; $f(-1 + \sqrt{6}) = (2 - 2\sqrt{6}) \cdot e^{\sqrt{6} - 1} \approx -12{,}35$ \quad f) $y = 3x + 5$ – an der $y$-Achse gespiegelt wechselt die Steigung das Vorzeichen, der $y$-Achsenabschnitt bleibt.}
\erg{45}{a) berührt – Funktionswert und Steigung stimmen überein. \quad b) $f(1) = 0 = g(1)$ und $f'(1) = 3 - 4 = -1$ = Steigung von $g$ – beide Bedingungen erfüllt: $g$ ist Tangente. \quad c) $f(x) = g(x) \Leftrightarrow 2x^2 - 4x + 2 = 2(x - 1)^2 = 0$, Stelle $x = 1$; $f(1) = g(1) = 2$, $f'(1) = g'(1) = 2$; gemeinsame Tangente $y = 2x$}
\erg{46}{a) $h(10) = 30$, $m = h'(10) = -\frac{1}{5} \cdot 10 = -2$; die Fortsetzung ist die Tangente: $y = -2(x - 10) + 30 = -2x + 50$ \quad b) $f(0) = -16$, $f'(0) = 24$, $t(x) = 24x - 16$; $f(x) - t(x) = x^3 - 9x^2 = x^2(x - 9) = 0$: $x = 0$ ist die Berührstelle, weiterer Punkt bei $x = 9$: $(9 | 200)$ \quad c) Berührpunkt $(2 | 3)$, Tangente $y = 2x - 1$ (ebenso möglich: Berührpunkt $(-2 | 3)$, $y = -2x - 1$) \quad d) Wendetangente: $f'(x) = \frac{3}{8}x^2 - \frac{3}{2}x$, $f'(2) = -1{,}5$; flachere fallende Geraden schneiden den Graphen noch zweimal, die Wendetangente und alle steileren nicht: $m \le -1{,}5$ (Rand eingeschlossen). \quad e) Wendetangente: $f'(x) = -x^2 + 2x$, $f'(1) = 1$; flachere steigende Geraden schneiden den Graphen noch zweimal, die Wendetangente und alle steileren nicht: $m \ge 1$ (Rand eingeschlossen).}

46c)

\begin{ksys}[xmin=-4,xmax=4,ymin=-2,ymax=6,klein] \funktion{0.5*\x^2+1}{f} \gerade{2}{-1}{t} \gerade{-2}{-1}{t_2} \end{ksys}

\erg{47}{a) ein Winkel – über den Tangens in die Steigung übersetzen, dann $f'(x) = m$ lösen. \quad b) $\alpha = \mathrm{tan}^{-1}(2) \approx 63{,}4^\circ$ \quad c) $f'(x) = x^2 - 2x$, $f'(3) = 3$; $\alpha = \mathrm{tan}^{-1}(3) \approx 71{,}6^\circ$}
\erg{48}{a) $f(1) = \frac{3}{e} \approx 1{,}10$, $m = f'(1) = -\frac{1}{e} \approx -0{,}37$; Achsenabschnitt aus dem Punkt (nicht $f(1)$): $t(x) = -\frac{1}{e}x + \frac{4}{e} \approx -0{,}37x + 1{,}47$; Schnittwinkel mit der $x$-Achse $\mathrm{tan}^{-1}(\frac{1}{e}) \approx 20{,}2^\circ$ \quad b) $f(2) = 2$, $m = f'(2) = \frac{3}{8} \cdot 4 = 1{,}5$; $t(x) = 1{,}5x - 1$; $\alpha = \mathrm{tan}^{-1}(1{,}5) \approx 56{,}3^\circ$ \quad c) Schnittstelle: $e^{2x} = 2$, $x = \frac{1}{2}\mathrm{ln}\,2 \approx 0{,}35$; $s'(x) = 2e^{2x} = 2 \cdot 2 = 4$; Winkel zur waagerechten Geraden $\mathrm{tan}^{-1}(4) \approx 76{,}0^\circ$ (nicht der Winkel zur $y$-Achse) \quad d) $f'(0) = 0{,}5$; Steigungswinkel $\mathrm{tan}^{-1}(0{,}5) \approx 26{,}6^\circ$ und $\mathrm{tan}^{-1}(2) \approx 63{,}4^\circ$; Schnittwinkel $\approx 36{,}9^\circ$ (nicht $\mathrm{tan}^{-1}(2 - 0{,}5)$) \quad e) $f'(3) = 9$, $f''(3) = 6$; $|\mathrm{tan}^{-1}(9) - \mathrm{tan}^{-1}(6)| \approx 83{,}7^\circ - 80{,}5^\circ \approx 3{,}1^\circ$ \quad f) $f'(2) = 0{,}4$, $f'(16) = -0{,}3$; Steigungswinkel $\approx 21{,}8^\circ$ und $\approx -16{,}7^\circ$; Schnittwinkel $\approx 38{,}5^\circ$ – mehr als $30^\circ$, die Vorgabe wird nicht eingehalten. \quad g) Gefälle heißt fallend: $f'(x) = \mathrm{tan}(-30^\circ) \approx -0{,}58$ (nicht $+0{,}58$ und nicht $f(x) = -0{,}58$); mit dem Rechner $x \approx 1{,}26$, Punkt $Q(1{,}26 | 0{,}16)$ \quad h) $f'(3) = -6$; Steigungswinkel von $f$ in $S$: $\mathrm{tan}^{-1}(6) \approx 80{,}5^\circ$; $g$ liegt spiegelbildlich, der Winkel an der Spitze ist doppelt so groß: $\approx 161{,}1^\circ$ – die Vorgabe ist eingehalten (nicht den einfachen Winkel mit $155^\circ$ vergleichen). \quad i) $f'(4) = -1$; Steigungswinkel von $f$ in $S$: $\mathrm{tan}^{-1}(1) = 45^\circ$; der Winkel an der Spitze ist doppelt so groß: $2 \cdot 45^\circ = 90^\circ$ – kleiner als $100^\circ$, die Vorgabe ist verfehlt. \quad j) $f'(2) = 3 - 1 = 2$, Steigungswinkel $\approx 63{,}4^\circ$; $g$: $\mathrm{tan}^{-1}(\frac{1}{3}) \approx 18{,}4^\circ$; Schnittwinkel $\alpha \approx 45^\circ$ (genau, denn $\mathrm{tan}\,\alpha = (2 - \frac{1}{3}) : (1 + \frac{2}{3}) = 1$); $\mathrm{sin}\,\alpha = \mathrm{cos}\,\alpha$ gilt genau für diesen Winkel – die Aussage ist wahr. \quad k) $f'(2) = 12 - 4 = 8$, Steigungswinkel $\approx 82{,}9^\circ$; $g$: $\approx 63{,}4^\circ$; Schnittwinkel $\alpha \approx 19{,}4^\circ$; $\mathrm{sin}\,\alpha = \mathrm{cos}\,\alpha$ gälte nur für $45^\circ$ – die Aussage ist falsch.}
\erg{49}{a) Umfang – alle drei Seiten, die Hypotenuse mit dem Satz des Pythagoras. \quad b) Nullstelle $x = 2$, $y$-Achsenabschnitt $g(0) = 6$; $A = \frac{1}{2} \cdot 2 \cdot 6 = 6$ \quad c) $f'(4) = -2$, $t(x) = -2x + 8$; Katheten $4$ und $8$; $A = \frac{1}{2} \cdot 4 \cdot 8 = 16$; Hypotenuse $\sqrt{4^2 + 8^2} = \sqrt{80} \approx 8{,}94$; $U \approx 4 + 8 + 8{,}94 = 20{,}94$}
\erg{50}{a) $f'(x) = (1 - x) \cdot e^x$, $f'(0) = 1$; $t(x) = x + 2$; Dreieck $A = \frac{1}{2} \cdot 2 \cdot 2 = 2$ FE; $1$ FE $= 25$ m², also $A = 50$ m²; eingespart $58 - 50 = 8$ m² (Flächenmaßstab nicht vergessen) \quad b) $f(0) = 3$, $f'(0) = -2$, $t(x) = -2x + 3$; Ecken $(0 | 0)$, $(1{,}5 | 0)$, $(0 | 3)$; Hypotenuse $\sqrt{1{,}5^2 + 3^2} \approx 3{,}35$; $U \approx 1{,}5 + 3 + 3{,}35 = 7{,}85$; gleicher Abstand: Mitte der Hypotenuse $(0{,}75 | 1{,}5)$ (Thaleskreis, nicht der Schwerpunkt) \quad c) $f'(x) = -\mathrm{cos}(x)$, $f'(2k\pi) = -1$ und $f(2k\pi) = 3$; Tangente $y = -x + 2k\pi + 3$; beide Achsenabschnitte sind $2k\pi + 3$ – für jedes $k$ gleich lang, also gleichschenklig (allgemein, nicht nur für ein $k$). \quad d) $h'(0) = -2$, $t(x) = -2x$; Dreieck unterhalb der $x$-Achse: $A = \frac{1}{2} \cdot b \cdot |t(b)| = b^2 = 49$, also $b = 7$ (Fläche positiv ansetzen)}
